% TOP_PROBLEM: {"id": 20000926, "title": "A saturated-list kernel for one-point extensions of four-color Ramsey colorings", "category_id": 2, "category": {"id": 2, "name": "combinatorics", "display_name": "Combinatorics", "description": "Counting problems, graph theory, discrete structures.", "slug": "combinatorics", "order_index": 2, "created_at": "2026-07-31T15:26:25.671Z"}, "status": "partially_solved", "problem_number": "AIM-COMBINATORICS-0051", "set_id": 14}
% FIRST_POSTED: 2026-10-01
% ENTRY_KIND: statement_only
% UnsolvedMath ID 20000926; source catalogue code AIM-COMBINATORICS-0051
% !TeX program = lualatex
% Definition and sources only; this file is not an LLM solution attempt.
\documentclass[11pt,a4paper]{article}
\usepackage[margin=25mm]{geometry}
\usepackage{fontspec}
\setmainfont{lmroman10-regular.otf}[BoldFont=lmroman10-bold.otf,
 ItalicFont=lmroman10-italic.otf,BoldItalicFont=lmroman10-bolditalic.otf]
\usepackage{amsmath,amssymb,mathtools}
\usepackage{unicode-math}
\setmathfont{latinmodern-math.otf}
\AtBeginDocument{\let\setminus\smallsetminus}
\usepackage{xurl,hyperref}
\hypersetup{colorlinks=true,urlcolor=blue}
\setcounter{secnumdepth}{0}
\setlength{\parindent}{0pt}
\setlength{\parskip}{5pt}
\setlength{\emergencystretch}{3em}
\begin{document}
\section*{A saturated-list kernel for one-point extensions of four-color Ramsey colorings}\label{20000926}
{\small UnsolvedMath ID 20000926 · AIM-COMBINATORICS-0051 · Combinatorics · AIM Workshop Problem Lists}
\subsection{Definitions and mathematical statement}
For an integer $N\ge1$, let $K_N$ be the complete graph on $[N]$,
with one undirected edge for each unordered pair of distinct vertices.
A four-coloring is a function
$c:\binom{[N]}2\to\{1,2,3,4\}$; colors are allowed to go unused.
A monochromatic triangle is a set of distinct vertices $x,y,z$ such
that $c(xy)=c(xz)=c(yz)$.
Define the four-color triangle Ramsey number by
\[
 R_4(3)=\min\{N:\text{every four-coloring of }K_N
                           \text{ has a monochromatic triangle}\}.
\]
This is the number also denoted $r(3,3,3,3)$.
The problem asks for its exact value and specifically asks whether
$R_4(3)=51$.

An exact determination $R_4(3)=M$ has two complementary requirements:
a four-coloring of $K_{M-1}$ without monochromatic triangles, and a
proof that every four-coloring of $K_M$ has one. A proposed coloring
must specify the color of every edge, including edges between any
parts used in its construction. The upper-bound assertion applies to
all colorings, without symmetry or algebraic restrictions.
Equivalently, determine the largest $N$ for which the edges of $K_N$
can be partitioned into four triangle-free graphs on the same vertex
set. The word ``triangle-free'' concerns ordinary triangles in a color
class; triangles whose edges have several different colors are allowed.
\subsection{Short English statement}
Find the exact four-color triangle Ramsey number; in particular, decide whether it equals 51.
\subsection{Sources}
\begin{itemize}
\item[S1] ULAM.AI and the UnsolvedMath Contributors, supplied problems.json, record 20000926 (AIM-COMBINATORICS-0051), AIM Workshop Problem Lists. Catalogue identity and source transcription; CC BY 4.0.
\url{https://huggingface.co/datasets/ulamai/UnsolvedMath}.
\item[S2] American Institute of Mathematics, Graph Ramsey theory, item 4. Original workshop question underlying AIM-COMBINATORICS-0051.
\url{https://www.overleaf.com/read/mnvcscjjysvg}.
\end{itemize}
\end{document}
